\documentclass[10pt,a4paper]{article}
\usepackage[utf8]{inputenc}
\usepackage[T1]{fontenc}
\usepackage[ngerman]{babel}
\usepackage[margin=1.7cm,top=1.5cm,bottom=1.7cm]{geometry}
\usepackage{titlesec}
\titleformat*{\section}{\normalfont\large\bfseries}
\titleformat*{\subsection}{\normalfont\normalsize\bfseries}
\titlespacing*{\section}{0pt}{8pt}{3pt}
\titlespacing*{\subsection}{0pt}{6pt}{2pt}
\usepackage{amsmath,amssymb}
\usepackage{listings}
\usepackage{xcolor}
\usepackage{tikz}
\usepackage{fancyhdr}
\usepackage{parskip}
\usepackage{needspace}

\newcommand{\mcbox}{\raisebox{-1pt}{\fbox{\rule{0pt}{7pt}\rule{7pt}{0pt}}}\hspace{6pt}}
\newcommand{\antwortlinien}[1]{%
  \par\vspace{2pt}%
  \newcount\zeilen \zeilen=#1
  \loop
    \noindent\rule{\linewidth}{0.4pt}\par\vspace{11pt}%
    \advance\zeilen by -1
  \ifnum\zeilen>0 \repeat
}
\lstset{language=Python,basicstyle=\ttfamily\footnotesize,keywordstyle=\bfseries,
  commentstyle=\itshape\color{gray},showstringspaces=false,frame=single,
  framerule=0.3pt,xleftmargin=8pt,framexleftmargin=6pt,aboveskip=6pt,belowskip=6pt}

\pagestyle{fancy}
\fancyhf{}
\fancyhead[L]{\small MDT5/2 -- Session 2: Abstandsbasierte Verfahren}
\fancyhead[R]{\small Übungsaufgaben zur Klausurvorbereitung}
\fancyfoot[C]{\small Seite \thepage}
\renewcommand{\headrulewidth}{0.4pt}

\begin{document}

\begin{center}
  {\Large\bfseries Übungsaufgaben Session 2}\\[4pt]
  {\large Kapitel 4: k-NN, Isolation Forest -- Praktikum 02}\\[4pt]
  \small Stil der Übungssammlung MDT5/2 -- generierte Aufgaben, keine Originale
\end{center}

\vspace{4pt}

\section*{Aufgabe 1 -- Multiple Choice mit Begründung}
\emph{Markieren Sie jeweils die einzige richtige von drei Aussagen und begründen Sie Ihre Entscheidung.}

\subsection*{1.1 \ Für k-Nearest-Neighbor zur Novelty Detection gilt:}
\mcbox Das Training ist aufwendig, dafür ist die Anwendung auf neue Punkte sehr schnell.

\mcbox Im Training werden lediglich alle normalen Trainingspunkte gespeichert; als Anomalie-Score
kann z.\,B. der mittlere Abstand eines neuen Punkts zu seinen $k$ nächsten Trainingspunkten dienen.

\mcbox Das Verfahren funktioniert umso besser, je höher die Anzahl der Dimensionen ist.

Begründung:
\antwortlinien{2}

\subsection*{1.2 \ Für Isolation Forests gilt:}
\mcbox Anomalien erzeugen in den Isolation Trees im Mittel längere Pfade als Normalpunkte.

\mcbox Eine Normalisierung der Eingabedaten ist für Isolation Forests essentiell.

\mcbox Die Splits erfolgen an zufälligen Merkmalen mit zufälligen Schwellwerten; Anomalien
lassen sich dadurch im Mittel mit weniger Splits isolieren.

Begründung:
\antwortlinien{2}

\section*{Aufgabe 2 -- k-NN-Score von Hand rechnen}

Die Normaldaten im Training sind die vier Punkte
$(0,0)$, $(0,2)$, $(2,0)$, $(2,2)$.
Als Anomalie-Score wird der \textbf{mittlere euklidische Abstand zu den $k=2$ nächsten
Trainingspunkten} verwendet.

a) Berechnen Sie den Score für den neuen Punkt $(1,1)$.
\antwortlinien{2}

b) Berechnen Sie den Score für den neuen Punkt $(6,1)$.
\antwortlinien{2}

c) Bei welchem Schwellwert würden Sie intuitiv die Grenze zwischen normal und Anomalie ziehen
und warum ist die Wahl des Schwellwerts generell problematisch ohne Validierungsdaten mit Anomalien?
\antwortlinien{3}

d) Nennen Sie zwei Schwächen von k-NN aus der Vorlesung.
\antwortlinien{2}

\section*{Aufgabe 3 -- Isolation-Forest-Score interpretieren}

Der normalisierte Anomalie-Score eines Isolation Forests ist
$s(\boldsymbol{x},n) = 2^{-\frac{E(h(\boldsymbol{x}))}{c(n)}}$,
wobei $E(h(\boldsymbol{x}))$ die mittlere Pfadlänge des Punkts über alle Bäume und
$c(n)$ die mittlere Pfadlänge einer erfolglosen Suche im Binärbaum ist.
Gegeben sei $c(n) = 10{,}2$.

a) Berechnen Sie $s$ für einen Punkt mit $E(h) = 3{,}4$ und interpretieren Sie das Ergebnis.
\antwortlinien{2}

b) Berechnen Sie $s$ für einen Punkt mit $E(h) = 10{,}2$ und interpretieren Sie das Ergebnis.
\antwortlinien{2}

c) Welche drei Grenzfälle für $s$ kennt die Vorlesung ($E(h)\to 0$, $E(h)\to c(n)$, $E(h)\to n-1$)
und was bedeuten sie?
\antwortlinien{3}

d) Warum wird pro Baum Subsampling verwendet (guter Standardwert $n_{sub}=256$),
und warum wird die maximale Baumtiefe auf $\log_2(n_{sub})$ begrenzt?
\antwortlinien{3}

\section*{Aufgabe 4 -- Verfahrenswahl: k-NN oder Isolation Forest?}

Ein Datensatz für Novelty Detection enthält zwei Merkmale mit sehr unterschiedlichen
Wertebereichen: $x_1 \in [0;\,1]$ (normierter Druck) und $x_2 \in [0;\,1000]$
(Temperatur-Rohwert).

a) Die Daten werden \emph{ohne} Normalisierung mit dem k-NN-Abstandsverfahren ausgewertet.
Welches Problem entsteht, und welches der beiden Merkmale dominiert den Score?
Nennen Sie die beiden Normalisierungsverfahren aus der Vorlesung, die das Problem beheben.
\antwortlinien{3}

b) Begründen Sie, warum ein Isolation Forest auch ohne Normalisierung funktioniert.
\antwortlinien{2}

c) Ein zweiter Datensatz hat mehrere hundert Merkmale und sehr viele Trainingspunkte.
Welches der beiden Verfahren ist hier besser geeignet? Nennen Sie zwei Gründe aus der Vorlesung.
\antwortlinien{3}

\section*{Aufgabe 5 -- Code-Verständnis}

\Needspace*{7cm}
Beurteilen Sie, ob der folgende Code für \textbf{Novelty Detection} (Training nur mit Normaldaten,
Anwendung auf neue Daten) geeignet ist, und begründen Sie:

\begin{lstlisting}
from sklearn.ensemble import IsolationForest

clf = IsolationForest(n_estimators=100,
                      max_samples=256,
                      contamination=0.01,
                      random_state=42)
clf.fit(normal_train_x)
predictions = clf.predict(test_x)   # +1 = normal, -1 = Anomalie
\end{lstlisting}

a) Erläutern Sie die Bedeutung der Parameter \texttt{n\_estimators}, \texttt{max\_samples}
und \texttt{contamination}.
\antwortlinien{3}

b) Ist der Ablauf (fit auf Normaldaten, predict auf Testdaten) für Novelty Detection korrekt?
\antwortlinien{2}

\vspace{10pt}
\begin{center}
\footnotesize\itshape
Nach Bearbeitung: Antworten an Claude schicken (Foto genügt) -- Korrektur mit Begründungs-Check.
\end{center}

\end{document}
